\documentclass[10pt,a4paper]{amsart}
\usepackage[margin=19mm]{geometry}
\usepackage{amsmath,amssymb,mathtools}
\usepackage[scr=rsfs,cal=euler]{mathalfa}
\usepackage{xfrac}
\usepackage[unicode,hidelinks,bookmarks=false]{hyperref}
\newcommand{\ie}{i.e.\ }
\newcommand{\cf}{cf.\ }
\newcommand{\resp}{respectively\ }
\newcommand{\Subsection}{\subsection}
\providecommand{\nobreakdash}{\nobreak\hbox{-}}
\setlength{\emergencystretch}{2em}
\multlinegap0pt
\usepackage{amssymb}
\usepackage{tikz}
\usepackage{tcolorbox}
\usepackage[shortlabels]{enumitem}
\newtheorem{lem}{Lemma}
\title{\Large{\bf{{Stochastic stability for weakly hyperbolic contracting Lorenz maps}}}}
\author{Haoyang Ji}
\date{Source: arXiv:2604.08006v1 (CC BY 4.0)}
\begin{document}
\maketitle

\noindent\emph{Source and licence.} \Large{\bf{{Stochastic stability for weakly hyperbolic contracting Lorenz maps}}}. Version arXiv:2604.08006v1 (CC BY 4.0), distributed by arXiv under the Creative Commons Attribution 4.0 International licence, http://creativecommons.org/licenses/by/4.0/, as recorded in the arXiv licence metadata for this exact version and shown on the abstract page https://arxiv.org/abs/2604.08006v1. Full licence notice: \texttt{licenses/arxiv-2604.08006-CC-BY-4.0.txt}.\par\smallskip

\noindent\textit{Editorial scope.} Counted unit: the article's lem lem (source0.tex lines 2682--2688). The article's complete original proof of it is delivered with it (source0.tex lines 2691--2775, one \texttt{proof} environment of 85 lines). No mathematics is written, repaired or re-derived: the statement and the proof are the article's own lines, in the article's own notation, and no retained line is substituted or re-flowed.  No further source-local context is needed: every cross-reference of every retained line resolves inside the two delivered spans.  Nothing was omitted from any retained span, so the package carries no omission record.  No retained line contains a citation, so the wrapper carries no bibliography and no bibliography entry.  No retained line contains a figure environment, an image-inclusion command or drawing code, so no image is delivered and the package declares no asset dependency.  Editorial formatting only, in addition: an AMS \texttt{amsart} preamble that restates the article's own declarations and macros used by the retained lines, namely the environments \texttt{lem} on separate counters, together with the standard packages those lines need.  The retained environments print in this document as Lemma 1 in order of appearance, whereas the article prints the same items with the numbering of its own full document.  The complete native source of the article as distributed in the arXiv e-print for this exact version is bundled unchanged as \texttt{sources/198129/source0.tex}.
\begin{lem}
Given $\theta > 0$ and $\gamma > 0$ there exists a constant $\tau > 0$ such that the following holds provided $\epsilon > 0$ is small enough. Let $(x, \omega) \in \tilde B(\epsilon) \times \Omega_{\epsilon}$, let $s \geq 1$ be an integer such that for each $0 \leq j < s$, 
\begin{equation}
s - j > \epsilon^{- \gamma} Q_{j}^{s-1} (x, \omega, \epsilon),
\end{equation}
then $\hat h_{\epsilon, \tau}^{\theta} (x, \omega) \leq s$.
\end{lem}

\begin{proof}

Let $\beta = \gamma/4$. Let $\epsilon_0 > 0$ be small such that Lemma 5.5 holds for all $\delta \in (0, \epsilon_0]$. In the following we may assume that $\epsilon \in (0, \epsilon_0 /e]$ small. Let $N$ be the maximal positive integer such that $e^{N-1} \epsilon \leq \epsilon_0$, then 
\[
2 \leq N \leq \log \frac{\epsilon_0}{\epsilon} + 1 < \epsilon^{- \beta} \mbox{ provided $\epsilon > 0$  is small enough}.
\]


Let $s_0 = s$, for $i = 1, 2, \cdots, N$, define 
\[
s_i = \max\{0 \leq j \leq s : f_{\omega}^j (x) \in \tilde B(e^{N-i} \epsilon)  \}.
\]
Then $s_N \leq s_{N-1} \leq \cdots \leq s_1 \leq s_0$. Define an integer $n \in \{0, 1, \cdots, N-1\}$ in the following way: $n = N$, if $s_0 - s_N < \epsilon^{-3 \beta}$; otherwise, let $n$ be the minimal integer in $ \{0, 1, \cdots, N-1\}$ such that $s_n - s_{n+1} \geq {\epsilon}^{-2 \beta}$. Such minimal $n$ exists since for otherwise $s_0 - s_N = \sum_{i=0}^{N-1} (s_n - s_{n+1}) < N e^{-2 \beta} \leq {\epsilon}^{-3 \beta}$, a contradiction. By minimality of $n$, we have  $s_0 - s_n < N \epsilon^{-2 \beta} \leq \epsilon^{-3 \beta}$. By (5.9), it follows that for $0 \leq j < s_n$,
\begin{equation}
s_n - j = (s - j) - (s - s_n) \geq \epsilon^{- \gamma} Q_{j}^{s-1} (x, \omega, \epsilon) - \epsilon^{-3 \beta} \geq \epsilon^{- 3\beta} Q_{j}^{s_n-1} (x, \omega, \epsilon),
\end{equation}
when $Q_{j}^{s_n-1} (x, \omega, \epsilon)> 0$. Note that when $Q_{j}^{s_n-1} (x, \omega, \epsilon)= 0$, the inequality holds trivally. 


If $n = N$, $f_{\omega}^{s_N}(x) \in \tilde B(\epsilon) $. We can argue as in the proof of Lemma 5.6 to show that $s_N$ is a $\theta$-good return time of $(x, \omega)$ into the region $\tilde B(\epsilon) \times \Omega_{\epsilon}$. Indeed, Let $s_N > s_{N+1} > \cdots > s_{N+N_0} = 0$ be all the integers such that $f_{\omega}^{s_{N+i}}(x) \in \tilde B(\epsilon), 0 \leq i \leq N_0$. let 
\[
B_i = \frac{Df_{\omega}^{S_{N+i}}(x)}{d(f_{\omega}^{S_{N+i}}(x), c)} \mbox{ and } \tilde B_i = \frac{Df_{\omega}^{S_{N+i}}(x)}{|\tilde B(\epsilon)|}.
\]
Similarly, we have
\[
A(x, \omega, s_N) \leq B_{N_0} + (1+\theta(\epsilon)) \sum_{i = 1}^{N_0 -1} B_i + \theta(\epsilon) \tilde B_0.
\]
Choose $K_0$ large enough, by Lemma 5.5 statement (1) and (5.10), for each $1 \leq i \leq N_0$,
\begin{align*}
\log \frac{\tilde B_0}{B_i}& \geq K_0 \Gamma_{s_{N+i}}^{s_N-1} (x, \omega, \epsilon) - Q_{s_{N+i}}^{s_N-1} (x, \omega, \epsilon) + \epsilon^{\beta}(s_N - s_{N+i})\\
&  \geq K_0 i  - \epsilon^{3 \beta} (s_N - s_{N+i}) + \epsilon^{\beta} (s_N - s_{N+i}) \geq K_0 i.
\end{align*}
The last inequality is enough to obtain the desired result.


If $n < N$, then $s_n > s_{n+1} \geq s_N$, so $f_{\omega}^{s_n} (x) \notin \tilde B(\epsilon)$. For each $0 \leq j \leq s_n$, let 
\[
A_j := \frac{D f_{\omega}^j(x)}{ d(f_{\omega}^j(x), c)}.
\]
Now we estimate $A_{s_n} / A(x, \omega, s_n)$ from below. Let $( s_n > ) s_N > s_{N+1} > \cdots > s_{N+N_0} = 0$ be all the integers such that $f_{\omega}^{s_{N+i}}(x) \in \tilde B(\epsilon), 0 \leq i \leq N_0$. 


For $n \leq k \leq N+N_0$, by Lemma 5.5 statement (2), we have
\[
\log \frac{A_{s_n}}{A_{s_k}} \geq \epsilon^{\beta} (s_n - s_k) + 2 \log \epsilon - Q_{s_k}^{s_n -1}(x, \omega, \epsilon).
\]
To apply Lemma 5.5, we set $K=1$, $(y, \overline \omega) = F^{s_k}(x, \omega)$, and $\delta = \epsilon$ in the case that $k \geq N$ and $\delta = e^{N-k} \epsilon$ for otherwise, the estimate keeps the same. Since $s_n - s_k \geq s_n - s_{n+1} \geq \epsilon^{-2 \beta}$, by (5.10) we have
\begin{align*}
\log \frac{A_{s_n}}{A_{s_k}} & \geq \epsilon^{\beta} (s_n - s_k) + 2 \log \epsilon - \epsilon^{3 \beta} (s_n - s_k) \geq \frac{\epsilon^{\beta}}{2} (s_n - s_k) \\
& \geq \frac{\epsilon^{\beta}}{2} (s_n - s_{n+1}) + \frac{\epsilon^{\beta}}{2} (s_{n+1} - s_k) \geq \frac{\epsilon^{-\beta}}{2} + \frac{\epsilon^{\beta}}{2} \max\{ k- N, 0 \}.
\end{align*}
Thus, 
\[
\frac{\sum_{i=n+1}^{N + N_0} A_{s_i}}{A_{s_n}} \leq \sum_{i=n+1}^{N + N_0} \frac{1}{e^{\frac{1}{2 \epsilon^{\beta}} +  \frac{\epsilon^{\beta}}{2} \max\{ i- N, 0 \}}}
\]
is very large. By Proposition 9, this implies 
\begin{align*}
A(x, \omega, s_{n+1} + 1)& = \sum_{i=S_{N+N_0}}^{s_{n+1}} \frac{D f_{\omega}^j(x)}{ d(f_{\omega}^j(x), c)} = \sum_{i = N+N_0}^{n+1} A_{s_i} + \sum_{i = N+N_0}^{n+2} \sum_{j = s_i+1}^{s_{i-1} -1} A_j \\
& \leq 2 \sum_{i = n+1}^{N + N_0} A_{s_i} \ll A_{s_n}. 
\end{align*}
To finish the proof, we distinguish two cases.


{\it Case 1.} $n > 0$. Then $f_{\omega}^{s_n} (x) \in \tilde B(e^{N-n} \epsilon) \subset \tilde B(\epsilon_0)$. By Proposition 9 again, 
\[
\sum_{j = s_{n+1} +1}^{s_n -1}  \frac{D f_{\omega}^j(x)}{ d(f_{\omega}^j(x), c)} \leq \theta(\epsilon_0) \frac{Df_{\omega}^{s_n}(x)}{ |\tilde B(\epsilon_0)|} \ll  \frac{D f_{\omega}^{s_n}(x)}{ d(f_{\omega}^{s_n}(x), c)} = A_{s_n},
\]
provided $\epsilon_0 > 0$ is small enough. This and the above inequality shows that $A(x, \omega, s_n) \ll A_{s_n}$. Note that since $f_{\omega}^{s_n} (x) \notin \tilde B(\epsilon)$,
\[
A_{s_n} = \frac{Df_{\omega}^{s_n}(x)}{|\tilde B(\epsilon)|} \frac{|\tilde B(\epsilon) |}{d(f_{\omega}^{s_n}(x), c)} \leq 2  \frac{Df_{\omega}^{s_n}(x)}{|\tilde B(\epsilon)|}.
\]
This implies $s_n$ is a $\theta$-good return time of $(x, \omega)$ into $\tilde B(\epsilon) \times \Omega_{\epsilon}$.


{\it Case 2.} $n =0$. Since $s_1 < s_0$ is well-defined, then $f_{\omega}^j(x) \notin \tilde B(\epsilon_0/e)$ for all $s_1 < j \leq s_0$. Note that in this case, $A_{s_0} \leq Df_{\omega}^{s_0} (x) / |\tilde B(\epsilon_0/e)|$. By Theorem 2 statement (2), $Df_{\omega}^{s_0} / Df_{\omega}^j (x)$ is exponentially large in $s_0 -j$, hence 
\[
\sum_{j = s_1 +1}^{s_0 -1}  \frac{D f_{\omega}^j(x)}{ d(f_{\omega}^j(x), c)} \leq C  \frac{D f_{\omega}^{s_0}(x)}{ |\tilde B(\epsilon_0/e)|},
\]
where $C$ depends only on $\ell$ and $\epsilon_0$. Combined with the above estimate, we have 
\[
A(x, \omega, s_0) = A(x, \omega, s_1 + 1) +  \sum_{j = s_1 +1}^{s_0 -1}  \frac{D f_{\omega}^j(x)}{ d(f_{\omega}^j(x), c)} \leq (C + \theta') \frac{Df_{\omega}^{s_0} (x)}{ |\tilde B(\epsilon_0/e)|},
\]
which implies that $s_0$ is a $\tau$-scale expansion time of $(x, \omega)$ for some constant $\tau > 0$ independent of $\epsilon$.

\end{proof}

\end{document}
